\documentclass[aps,preprint]{revtex4}%
\usepackage{amsfonts}
\usepackage{amsmath}
\usepackage{amssymb}
\usepackage{graphicx}%
\setcounter{MaxMatrixCols}{30}
%TCIDATA{OutputFilter=latex2.dll}
%TCIDATA{Version=5.50.0.2953}
%TCIDATA{CSTFile=revtex4.cst}
%TCIDATA{Created=Monday, August 09, 2004 09:20:59}
%TCIDATA{LastRevised=Saturday, July 28, 2007 09:51:31}
%TCIDATA{<META NAME="GraphicsSave" CONTENT="32">}
%TCIDATA{<META NAME="SaveForMode" CONTENT="1">}
%TCIDATA{BibliographyScheme=Manual}
%TCIDATA{<META NAME="DocumentShell" CONTENT="Articles\SW\REVTeX 4">}
%BeginMSIPreambleData
\providecommand{\U}[1]{\protect\rule{.1in}{.1in}}
%EndMSIPreambleData
\newtheorem{theorem}{Theorem}
\newtheorem{acknowledgement}[theorem]{Acknowledgement}
\newtheorem{algorithm}[theorem]{Algorithm}
\newtheorem{axiom}[theorem]{Axiom}
\newtheorem{claim}[theorem]{Claim}
\newtheorem{conclusion}[theorem]{Conclusion}
\newtheorem{condition}[theorem]{Condition}
\newtheorem{conjecture}[theorem]{Conjecture}
\newtheorem{corollary}[theorem]{Corollary}
\newtheorem{criterion}[theorem]{Criterion}
\newtheorem{definition}[theorem]{Definition}
\newtheorem{example}[theorem]{Example}
\newtheorem{exercise}[theorem]{Exercise}
\newtheorem{lemma}[theorem]{Lemma}
\newtheorem{notation}[theorem]{Notation}
\newtheorem{problem}[theorem]{Problem}
\newtheorem{proposition}[theorem]{Proposition}
\newtheorem{remark}[theorem]{Remark}
\newtheorem{solution}[theorem]{Solution}
\newtheorem{summary}[theorem]{Summary}
\newenvironment{proof}[1][Proof]{\noindent\textbf{#1.} }{\ \rule{0.5em}{0.5em}}
\begin{document}
\preprint{HEP/123-qed}
\title[Short title for running header]{Long title that appears on the title page}
\author{First Author}
\affiliation{My Institution}
\author{Second Author}
\affiliation{My Institution}
\author{Third Author}
\affiliation{Other Institution}
\keywords{one two three}
\pacs{PACS number}

\begin{abstract}
Shell document for REV\TeX{} 4.

\end{abstract}
\volumeyear{year}
\volumenumber{number}
\issuenumber{number}
\eid{identifier}
\date[Date text]{date}
\received[Received text]{date}

\revised[Revised text]{date}

\accepted[Accepted text]{date}

\published[Published text]{date}

\startpage{101}
\endpage{102}
\maketitle
\tableofcontents


\section{Variation of the Occupation Numbers\label{OccVar}}

In this section we show how by optimizing wrt geminal coefficients $\left\{
c_{j}\right\}  $ for a fixed value of $\left(  \mathbf{\mathbf{\psi},\theta
}\right)  $ one can determine an optimal value of $\mathbf{n}$ that satisfies
the positivity and normalization constraints.

The $2N$-electron, unnormalized, AGP state, $\left\vert g^{N}\right\rangle $,
can be expressed as
\begin{equation}
\left\vert g^{N}\right\rangle =\left(  G^{\dagger}\right)  ^{N}\left\vert
\phi\right\rangle ,
\end{equation}
where
\begin{equation}
G^{\dagger}=\sum_{1\leq j\leq s}c_{j}b_{j}^{\dagger}b_{j+s}^{\dagger},
\label{gcreat}%
\end{equation}
and the electron creation operators $\left\{  \left.  b_{j}^{\dagger
}\right\vert 1\leq j\leq2s\right\}  $ produce the CGSO's $\left\{  \left.
\left\vert \psi_{j}\right\rangle \right\vert 1\leq j\leq2s\right\}  $ by
acting on the zero particle vacuum vector i.e.
\begin{equation}
\left\vert \psi_{j}\right\rangle =b_{j}^{\dagger}\left\vert \phi\right\rangle
;\quad1\leq j\leq2s.
\end{equation}
For fixed $\left(  \mathbf{\mathbf{\psi},\theta}\right)  $ the geminal creator
in eq. $\left(  \ref{gcreat}\right)  $ is a function of the canonical
coefficients $\left\{  \left.  c_{j}\right\vert 1\leq j\leq s\right\}  ,$
which can be differentiated to give%
\begin{equation}
\frac{\partial G\left(  \mathbf{c}\right)  ^{\dagger N}}{\partial c_{j}%
}=Nb_{j}^{\dagger}b_{j+s}^{\dagger}G\left(  \mathbf{c}\right)  ^{\dagger N-1},
\end{equation}
which leads to
\begin{equation}
G\left(  \mathbf{c}\right)  ^{\dagger N}=\frac{1}{N}\sum_{1\leq j\leq s}%
c_{j}\frac{\partial G\left(  \mathbf{c}\right)  ^{\dagger N}}{\partial c_{j}}.
\label{euler}%
\end{equation}
The energy expression
\begin{equation}
\mathcal{E}\left(  \mathbf{c}\right)  =\frac{\mathcal{H}\left(  \mathbf{c}%
\right)  }{S_{N}\left(  \mathbf{c}\right)  }=\frac{\left\langle G\left(
\mathbf{c}\right)  ^{\dagger N}\phi|HG\left(  \mathbf{c}\right)  ^{\dagger
N}\phi\right\rangle }{\left\langle G\left(  \mathbf{c}\right)  ^{\dagger
N}\phi|G\left(  \mathbf{c}\right)  ^{\dagger N}\phi\right\rangle }%
\end{equation}
produces the derivative
\begin{equation}
\partial_{\mathbf{c}}^{1}\mathcal{E}\left(  \mathbf{c}\right)  =\frac{1}%
{S_{N}\left(  \mathbf{c}\right)  }\left\{  \partial_{\mathbf{c}}%
^{1}\mathcal{H}\left(  \mathbf{c}\right)  -\mathcal{E}\left(  \mathbf{c}%
\right)  \partial_{\mathbf{c}}^{1}S_{N}\left(  \mathbf{c}\right)  \right\}  ,
\end{equation}
where
\begin{equation}
\partial_{c_{j}}^{1}\mathcal{H}\left(  \mathbf{c}\right)  =N\left\{
\left\langle b_{j}^{\dagger}b_{j+s}^{\dagger}G\left(  \mathbf{c}\right)
^{\dagger N-1}\phi|HG\left(  \mathbf{c}\right)  ^{\dagger N}\phi\right\rangle
+\left\langle G\left(  \mathbf{c}\right)  ^{\dagger N}\phi|Hb_{j}^{\dagger
}b_{j+s}^{\dagger}G\left(  \mathbf{c}\right)  ^{\dagger N-1}\phi\right\rangle
\right\}
\end{equation}
and
\begin{equation}
\partial_{c_{j}}^{1}S_{N}\left(  \mathbf{c}\right)  =N\left\{  \left\langle
b_{j}^{\dagger}b_{j+s}^{\dagger}G\left(  \mathbf{c}\right)  ^{\dagger N-1}%
\phi|G\left(  \mathbf{c}\right)  ^{\dagger N}\phi\right\rangle +\left\langle
G\left(  \mathbf{c}\right)  ^{\dagger N}\phi|b_{j}^{\dagger}b_{j+s}^{\dagger
}G\left(  \mathbf{c}\right)  ^{\dagger N-1}\phi\right\rangle \right\}  .
\end{equation}
Using eq. $\left(  \ref{euler}\right)  $ one obtains
\begin{equation}
\frac{1}{2N}\sum_{1\leq j\leq s}c_{j}\frac{\partial\mathcal{E}\left(
\mathbf{c}\right)  }{\partial c_{j}}=\mathbf{c}^{t}\partial_{\mathbf{c}}%
^{1}\mathcal{E}\left(  \mathbf{c}\right)  =\frac{1}{2NS_{N}\left(
\mathbf{c}\right)  }\left\{  \mathcal{H}\left(  \mathbf{c}\right)
-\mathcal{E}\left(  \mathbf{c}\right)  S_{N}\left(  \mathbf{c}\right)
\right\}  =0,
\end{equation}
i.e. the vector $\mathbf{c}$ is always perpendicular to the vector
$\partial_{\mathbf{c}}^{1}\mathcal{E}\left(  \mathbf{c}\right)  $.

The second derivatives $\left\{  \partial_{c_{j}c_{k}}^{2}\mathcal{E}\left(
\mathbf{c}\right)  \right\}  $ are given by
\begin{equation}
\partial_{\mathbf{c}}^{2}\mathcal{E}\left(  \mathbf{c}\right)  =\frac{1}%
{S_{N}\left(  \mathbf{c}\right)  }\left\{  \partial_{\mathbf{c}}%
^{2}\mathcal{H}\left(  \mathbf{c}\right)  -\mathcal{E}\left(  \mathbf{c}%
\right)  \partial_{\mathbf{c}}^{2}S_{N}\left(  \mathbf{c}\right)
-\partial_{\mathbf{c}}^{1}\mathcal{E}\left(  \mathbf{c}\right)  \otimes
\partial_{\mathbf{c}}^{1}S_{N}\left(  \mathbf{c}\right)  -\partial
_{\mathbf{c}}^{1}S_{N}\left(  \mathbf{c}\right)  \otimes\partial_{\mathbf{c}%
}^{1}\mathcal{E}\left(  \mathbf{c}\right)  \right\}  , \label{Hess}%
\end{equation}
where
\begin{align}
\partial_{c_{j}c_{k}}^{2}\mathcal{H}\left(  \mathbf{c}\right)   &
=N^{2}\left\{  \left\langle b_{j}^{\dagger}b_{j+s}^{\dagger}G\left(
\mathbf{c}\right)  ^{\dagger N-1}\phi|Hb_{k}^{\dagger}b_{k+s}^{\dagger
}G\left(  \mathbf{c}\right)  ^{\dagger N-1}\right\rangle +\left\langle
b_{k}^{\dagger}b_{k+s}^{\dagger}G\left(  \mathbf{c}\right)  ^{\dagger N-1}%
\phi|Hb_{j}^{\dagger}b_{j+s}^{\dagger}G\left(  \mathbf{c}\right)  ^{\dagger
N-1}\right\rangle \right\} \nonumber\\
&  +N\left(  N-1\right)  \left\{  \left\langle G\left(  \mathbf{c}\right)
^{\dagger N}\phi|Hb_{j}^{\dagger}b_{j+s}^{\dagger}b_{k}^{\dagger}%
b_{k+s}^{\dagger}G\left(  \mathbf{c}\right)  ^{\dagger N-2}\phi\right\rangle
\right. \nonumber\\
&  \qquad\qquad\qquad\qquad\qquad\qquad\qquad\qquad\qquad\left.  +\left\langle
b_{j}^{\dagger}b_{j+s}^{\dagger}b_{k}^{\dagger}b_{k+s}^{\dagger}G\left(
\mathbf{c}\right)  ^{\dagger N-2}\phi|HG\left(  \mathbf{c}\right)  ^{\dagger
}\phi\right\rangle \right\}
\end{align}
and
\begin{align}
\partial_{c_{j}c_{k}}^{2}S_{N}\left(  \mathbf{c}\right)   &  =N^{2}\left\{
\left\langle b_{j}^{\dagger}b_{j+s}^{\dagger}G\left(  \mathbf{c}\right)
^{\dagger N-1}\phi|b_{k}^{\dagger}b_{k+s}^{\dagger}G\left(  \mathbf{c}\right)
^{\dagger N-1}\right\rangle +\left\langle b_{k}^{\dagger}b_{k+s}^{\dagger
}G\left(  \mathbf{c}\right)  ^{\dagger N-1}\phi|b_{j}^{\dagger}b_{j+s}%
^{\dagger}G\left(  \mathbf{c}\right)  ^{\dagger N-1}\right\rangle \right\}
\nonumber\\
&  +N\left(  N-1\right)  \left\{  \left\langle G\left(  \mathbf{c}\right)
^{\dagger N}\phi|b_{j}^{\dagger}b_{j+s}^{\dagger}b_{k}^{\dagger}%
b_{k+s}^{\dagger}G\left(  \mathbf{c}\right)  ^{\dagger N-2}\phi\right\rangle
\right. \nonumber\\
&  \qquad\qquad\qquad\qquad\qquad\qquad\qquad\qquad\qquad\left.  +\left\langle
b_{j}^{\dagger}b_{j+s}^{\dagger}b_{k}^{\dagger}b_{k+s}^{\dagger}G\left(
\mathbf{c}\right)  ^{\dagger N-2}\phi|G\left(  \mathbf{c}\right)  ^{\dagger
}\phi\right\rangle \right\}  .
\end{align}
Observing
\begin{align}
&  \sum_{1\leq k\leq s}\left\{  \left\langle b_{j}^{\dagger}b_{j+s}^{\dagger
}G\left(  \mathbf{c}\right)  ^{\dagger N-1}\phi|Hb_{k}^{\dagger}%
b_{k+s}^{\dagger}G\left(  \mathbf{c}\right)  ^{\dagger N-1}\right\rangle
+\left\langle b_{k}^{\dagger}b_{k+s}^{\dagger}G\left(  \mathbf{c}\right)
^{\dagger N-1}\phi|Hb_{j}^{\dagger}b_{j+s}^{\dagger}G\left(  \mathbf{c}%
\right)  ^{\dagger N-1}\right\rangle \right\}  c_{k}\nonumber\\
&  =\left\langle b_{j}^{\dagger}b_{j+s}^{\dagger}G\left(  \mathbf{c}\right)
^{\dagger N-1}\phi|HG\left(  \mathbf{c}\right)  ^{\dagger N}\right\rangle
+\left\langle G\left(  \mathbf{c}\right)  ^{\dagger N}\phi|Hb_{j}^{\dagger
}b_{j+s}^{\dagger}G\left(  \mathbf{c}\right)  ^{\dagger N-1}\right\rangle
=\frac{1}{N}\partial_{c_{j}}^{1}\mathcal{H}\left(  \mathbf{c}\right)  ,
\end{align}
and
\begin{align}
&  \sum_{1\leq k\leq s}\left\{  \left\langle G\left(  \mathbf{c}\right)
^{\dagger N}\phi|Hb_{j}^{\dagger}b_{j+s}^{\dagger}b_{k}^{\dagger}%
b_{k+s}^{\dagger}G\left(  \mathbf{c}\right)  ^{\dagger N-2}\phi\right\rangle
+\left\langle b_{j}^{\dagger}b_{j+s}^{\dagger}b_{k}^{\dagger}b_{k+s}^{\dagger
}G\left(  \mathbf{c}\right)  ^{\dagger N-2}\phi|HG\left(  \mathbf{c}\right)
^{\dagger}\phi\right\rangle \right\}  c_{k}\nonumber\\
&  =\left\langle G\left(  \mathbf{c}\right)  ^{\dagger N}\phi|Hb_{j}^{\dagger
}b_{j+s}^{\dagger}G\left(  \mathbf{c}\right)  ^{\dagger N-1}\phi\right\rangle
+\left\langle b_{j}^{\dagger}b_{j+s}^{\dagger}G\left(  \mathbf{c}\right)
^{\dagger N-1}\phi|HG\left(  \mathbf{c}\right)  ^{\dagger N}\phi\right\rangle
=\frac{1}{N}\partial_{c_{j}}^{1}\mathcal{H}\left(  \mathbf{c}\right)
\end{align}
one obtains
\begin{equation}
\sum_{1\leq k\leq s}\partial_{c_{j}c_{k}}^{2}\mathcal{H}\left(  \mathbf{c}%
\right)  c_{k}=\left(  2N-1\right)  \partial_{c_{j}}^{1}\mathcal{H}\left(
\mathbf{c}\right)  .
\end{equation}
The normalization term similarly gives
\begin{align}
&  \sum_{1\leq k\leq s}\left\{  \left\langle b_{j}^{\dagger}b_{j+s}^{\dagger
}G\left(  \mathbf{c}\right)  ^{\dagger N-1}\phi|b_{k}^{\dagger}b_{k+s}%
^{\dagger}G\left(  \mathbf{c}\right)  ^{\dagger N-1}\right\rangle
+\left\langle b_{k}^{\dagger}b_{k+s}^{\dagger}G\left(  \mathbf{c}\right)
^{\dagger N-1}\phi|b_{j}^{\dagger}b_{j+s}^{\dagger}G\left(  \mathbf{c}\right)
^{\dagger N-1}\right\rangle \right\}  c_{k}\nonumber\\
&  =\left\langle b_{j}^{\dagger}b_{j+s}^{\dagger}G\left(  \mathbf{c}\right)
^{\dagger N-1}\phi|G\left(  \mathbf{c}\right)  ^{\dagger N}\right\rangle
+\left\langle G\left(  \mathbf{c}\right)  ^{\dagger N}\phi|b_{j}^{\dagger
}b_{j+s}^{\dagger}G\left(  \mathbf{c}\right)  ^{\dagger N-1}\right\rangle
=\frac{1}{N}\partial_{c_{j}}^{1}S_{N}\left(  \mathbf{c}\right)  ,
\end{align}
and
\begin{align}
&  \sum_{1\leq k\leq s}\left\{  \left\langle G\left(  \mathbf{c}\right)
^{\dagger N}\phi|b_{j}^{\dagger}b_{j+s}^{\dagger}b_{k}^{\dagger}%
b_{k+s}^{\dagger}G\left(  \mathbf{c}\right)  ^{\dagger N-2}\phi\right\rangle
+\left\langle b_{j}^{\dagger}b_{j+s}^{\dagger}b_{k}^{\dagger}b_{k+s}^{\dagger
}G\left(  \mathbf{c}\right)  ^{\dagger N-2}\phi|G\left(  \mathbf{c}\right)
^{\dagger}\phi\right\rangle \right\}  c_{k}\nonumber\\
&  =\left\langle G\left(  \mathbf{c}\right)  ^{\dagger N}\phi|b_{j}^{\dagger
}b_{j+s}^{\dagger}G\left(  \mathbf{c}\right)  ^{\dagger N-1}\phi\right\rangle
+\left\langle b_{j}^{\dagger}b_{j+s}^{\dagger}G\left(  \mathbf{c}\right)
^{\dagger N-1}\phi|G\left(  \mathbf{c}\right)  ^{\dagger N}\phi\right\rangle
=\frac{1}{N}\partial_{c_{j}}^{1}S_{N}\left(  \mathbf{c}\right)  .
\end{align}
producing
\begin{equation}
\sum_{1\leq k\leq s}\partial_{c_{j}c_{k}}^{2}S_{N}\left(  \mathbf{c}\right)
c_{k}=\left(  2N-1\right)  \partial_{c_{j}}^{1}S_{N}\left(  \mathbf{c}\right)
.
\end{equation}
Thus
\begin{equation}
\sum_{1\leq k\leq s}\left\{  \partial_{c_{j}c_{k}}^{2}\mathcal{H}\left(
\mathbf{c}\right)  -\mathcal{E}\left(  \mathbf{c}\right)  \partial_{c_{j}%
c_{k}}^{2}S_{N}\left(  \mathbf{c}\right)  \right\}  c_{k}=\left(  2N-1\right)
\left\{  \partial_{c_{j}}^{1}\mathcal{H}\left(  \mathbf{c}\right)
-\mathcal{E}\left(  \mathbf{c}\right)  \partial_{c_{j}}^{1}S_{N}\left(
\mathbf{c}\right)  \right\}  ,
\end{equation}
which can be written in matrix notation as
\begin{equation}
L\left(  \mathbf{c}\right)  \mathbf{c}\equiv\left[  \partial_{\mathbf{c}}%
^{2}\mathcal{H}\left(  \mathbf{c}\right)  -\mathcal{E}\left(  \mathbf{c}%
\right)  \partial_{\mathbf{c}}^{2}S_{N}\left(  \mathbf{c}\right)  \right]
\mathbf{c=}\left(  2N-1\right)  \partial_{\mathbf{c}}^{1}\mathcal{E}\left(
\mathbf{c}\right)  , \label{gradrel}%
\end{equation}
which clearly shows that
\begin{equation}
\mathbf{c}^{t}L\left(  \mathbf{c}\right)  \mathbf{c=\mathbf{c}}^{t}%
\mathbf{\partial_{\mathbf{c}}^{2}\mathcal{E}\left(  \mathbf{c}\right)
\mathbf{c}}=\mathbf{0.} \label{cd2c}%
\end{equation}


At stationary points, $\mathbf{c}_{\#}$, one can observe from eq. $\left(
\ref{Hess}\right)  $ that
\begin{equation}
\mathbf{\partial_{\mathbf{c}}^{2}\mathcal{E}\left(  \mathbf{c}_{\#}\right)
=}\frac{1}{S_{N}\left(  \mathbf{c}_{\#}\right)  }L\left(  \mathbf{c}%
_{\#}\right)
\end{equation}
and if $\mathbf{c}_{\#}$ is a local minimum then
\begin{equation}
L\left(  \mathbf{c}_{\#}\right)  \geq\mathbf{0.}%
\end{equation}
Hence at a local minimum eq. $\left(  \ref{cd2c}\right)  $ shows that
$\mathbf{c}_{\#}$ is an eigenvector corresponding to the lowest eigenvalue of
$L\left(  \mathbf{c}_{\#}\right)  $ i.e.%
\begin{equation}
L\left(  \mathbf{c}_{\#}\right)  \mathbf{c}_{\#}=0. \label{loweig}%
\end{equation}


If $\mathbf{c}_{\#}$ is the lowest eigenvalue of $L\left(  \mathbf{c}%
_{\#}\right)  $ and $L\left(  \mathbf{c}_{\#}\right)  \geq0$ then as
\begin{equation}
L\left(  \mathbf{c}_{\#}\right)  =S_{N}\left(  \mathbf{c}_{\#}\right)
\partial_{\mathbf{c}}^{2}\mathcal{E}\left(  \mathbf{c}_{\#}\right)
+\partial_{\mathbf{c}}^{1}\mathcal{E}\left(  \mathbf{c}_{\#}\right)
\otimes\partial_{\mathbf{c}}^{1}S_{N}\left(  \mathbf{c}_{\#}\right)
+\partial_{\mathbf{c}}^{1}S_{N}\left(  \mathbf{c}_{\#}\right)  \otimes
\partial_{\mathbf{c}}^{1}\mathcal{E}\left(  \mathbf{c}_{\#}\right)
\end{equation}
one has
\begin{equation}
L\left(  \mathbf{c}_{\#}\right)  \mathbf{c}_{\#}=S_{N}\left(  \mathbf{c}%
_{\#}\right)  \partial_{\mathbf{c}}^{2}\mathcal{E}\left(  \mathbf{c}%
_{\#}\right)  \mathbf{c}_{\#}+2NS_{N}\left(  \mathbf{c}_{\#}\right)
\partial_{\mathbf{c}}^{1}\mathcal{E}\left(  \mathbf{c}_{\#}\right)  =0,
\end{equation}
which implies
\begin{equation}
\frac{1}{2N}\partial_{\mathbf{c}}^{2}\mathcal{E}\left(  \mathbf{c}%
_{\#}\right)  \mathbf{c}_{\#}+\partial_{\mathbf{c}}^{1}\mathcal{E}\left(
\mathbf{c}_{\#}\right)  =0.\label{stat1}%
\end{equation}
If the second order approximation $\mathcal{E}_{2}$
\begin{equation}
\mathcal{E}\left(  \mathbf{c}_{\#}+\mathbf{\delta}\right)  \simeq
\mathcal{E}_{2}\left(  \mathbf{\delta}\right)  =\mathcal{E}\left(
\mathbf{c}_{\#}\right)  +\partial_{\mathbf{c}}^{1}\mathcal{E}\left(
\mathbf{c}_{\#}\right)  ^{t}\mathbf{\delta+}\frac{1}{2}\mathbf{\delta}%
^{t}\partial_{\mathbf{c}}^{2}\mathcal{E}\left(  \mathbf{c}_{\#}\right)
\mathbf{\delta}%
\end{equation}
to $\mathcal{E}$ about the point $\mathbf{c}_{\#}$ is stationary at the point
$\mathbf{\delta}_{\#}$ then
\begin{equation}
\partial_{\mathbf{\delta}}\mathcal{E}_{2}\left(  \mathbf{\delta}_{\#}\right)
=\partial_{\mathbf{c}}^{1}\mathcal{E}\left(  \mathbf{c}_{\#}\right)
^{t}+\mathbf{\delta}^{t}\partial_{\mathbf{c}}^{2}\mathcal{E}\left(
\mathbf{c}_{\#}\right)  =0
\end{equation}
and hence
\begin{equation}
\partial_{\mathbf{c}}^{1}\mathcal{E}\left(  \mathbf{c}_{\#}\right)
+\partial_{\mathbf{c}}^{2}\mathcal{E}\left(  \mathbf{c}_{\#}\right)
\mathbf{\delta}=0\label{stat2}%
\end{equation}
Thus
\begin{equation}
\mathbf{\delta}_{\#}=\frac{1}{2N}\mathbf{c}_{\#}%
\end{equation}
is a solution to eq. $\left(  \ref{stat2}\right)  $ by observing eq. $\left(
\ref{stat1}\right)  .$ However as $\mathcal{E}$ is a homogenous function of
degree $0$
\begin{equation}
\mathcal{E}\left(  \lambda\mathbf{c}\right)  =\mathcal{E}\left(
\mathbf{c}\right)  \quad\forall\lambda\in\mathbb{R}%
\end{equation}
and
\begin{equation}
\mathcal{E}\left(  \mathbf{c}_{\#}+\mathbf{\delta}_{\#}\right)  =\mathcal{E}%
\left(  \mathbf{c}_{\#}+\frac{1}{2N}\mathbf{c}_{\#}\right)  =\mathcal{E}%
\left(  \mathbf{c}_{\#}\right)  .
\end{equation}
This means $\mathbf{\delta}_{\#}=\mathbf{0}$ is also a stationary point of
$\mathcal{E}_{2}\left(  \mathbf{\delta}_{\#}\right)  $ and hence eq. $\left(
\ref{stat2}\right)  $ shows that
\begin{equation}
\partial_{\mathbf{c}}^{1}\mathcal{E}\left(  \mathbf{c}_{\#}\right)
=\mathbf{0,}%
\end{equation}
i.e. $\mathbf{c}_{\#}$ is stationary point of $\mathcal{E}.$ At stationary
points of $\mathcal{E}$
\begin{equation}
\mathbf{\partial_{\mathbf{c}}^{2}\mathcal{E}\left(  \mathbf{c}_{\#}\right)
=}\frac{1}{S_{N}\left(  \mathbf{c}_{\#}\right)  }L\left(  \mathbf{c}%
_{\#}\right)
\end{equation}
and by assumption $L\left(  \mathbf{c}_{\#}\right)  \geq0$ thus one can
conclude that $\mathbf{c}_{\#}$ is a local minimum. We can summarize the
preceding by the theorem

\begin{theorem}
$\mathbf{c}_{\#}$ is a local minimum of $\mathcal{E}$ if and only if
\begin{equation}
L\left(  \mathbf{c}_{\#}\right)  \geq0\qquad\&\qquad L\left(  \mathbf{c}%
_{\#}\right)  \mathbf{c}_{\#}=0.
\end{equation}

\end{theorem}

An iterative minimization scheme can be based on this theorem by considering
the first order approximation
\begin{equation}
\mathcal{E}\left(  \mathbf{c}+\mathbf{\delta}\right)  \simeq\mathcal{E}%
_{1}\left(  \mathbf{\delta}\right)  =\mathcal{E}\left(  \mathbf{c}\right)
+\partial_{\mathbf{c}}^{1}\mathcal{E}\left(  \mathbf{c}\right)  ^{t}%
\mathbf{\delta} \label{o1}%
\end{equation}
and the eigenvalue equation
\begin{equation}
L\left(  \mathbf{c}\right)  \mathbf{x}_{j}=\lambda_{j}\mathbf{x}_{j}\mathbf{,}
\label{Lc}%
\end{equation}
where $\left\{  \left.  \mathbf{x}_{j}\right\vert 1\leq j\leq s\right\}  $ are
orthonormal eigenvectors. By taking the scalar product with $\mathbf{c}$ eq.
$\left(  \ref{Lc}\right)  $ becomes by noting eq. $\left(  \ref{gradrel}%
\right)  $
\begin{equation}
\mathbf{c}^{t}L\left(  \mathbf{c}\right)  \mathbf{x}_{j}=\left(  2N-1\right)
\partial_{\mathbf{c}}^{1}\mathcal{E}\left(  \mathbf{c}\right)  ^{t}%
\mathbf{x}_{j}=\lambda_{j}\mathbf{c}^{t}\mathbf{x}_{j}%
\end{equation}
As $\mathbf{c}$ is always perpendicular to $\partial_{\mathbf{c}}%
^{1}\mathcal{E}\left(  \mathbf{c}\right)  $ we obtain
\begin{equation}
\partial_{\mathbf{c}}^{1}\mathcal{E}\left(  \mathbf{c}\right)  ^{t}\left(
\mathbf{x}_{j}-\mathbf{c}\right)  =\frac{1}{\left(  2N-1\right)  }\lambda
_{j}\mathbf{c}^{t}\mathbf{x}_{j}%
\end{equation}
We can always choose an overall phase for $\mathbf{x}_{j}$ so that
$\mathbf{c}^{t}\mathbf{x}_{j}>0.$ If $\lambda_{L}<0$ we can then set
\begin{equation}
\mathbf{\delta}=\alpha\left(  \mathbf{x}_{L}-\mathbf{c}\right)
\end{equation}
in eq. $\left(  \ref{o1}\right)  $, where $\lambda_{L}$ is the lowest
eigenvalue of $L\left(  \mathbf{c}\right)  $ and find a small enough
$\alpha>0$ that decreases the energy around the point $\mathbf{c.}$ Then we
can compute a new point
\begin{equation}
\mathbf{c}_{new}=\alpha\left(  \mathbf{x}_{L}-\mathbf{c}\right)
\end{equation}
and continue iterating until the lowest eigenvalue is equal to zero.

After the optimal $\mathbf{c}_{\#}$ has been found we can then define an
optimal $\mathbf{n}_{\#}$ by
\begin{equation}
n_{\#j}=c_{\#j}^{2}\qquad1\leq j\leq s.
\end{equation}
As the values $\left\{  c_{\#j}\right\}  $ are optimal
\begin{equation}
\frac{\partial\mathbf{\mathcal{E}\left(  \mathbf{c}_{\#}\right)  }}{\partial
c_{j}}=\frac{\partial\mathbf{\mathcal{E}\left(  \mathbf{n}_{\#}\right)  }%
}{\partial n_{j}}\frac{\partial n_{j}}{\partial c_{j}}=2\frac{\partial
\mathbf{\mathcal{E}\left(  \mathbf{n}_{\#}\right)  }}{\partial n_{j}}%
=0\qquad1\leq j\leq s,
\end{equation}
showing that $\left\{  n_{\#j}\right\}  $ are also optimal and satisfy the
positivity and normalization constraints. The phase information contained in
the signs of $\left\{  c_{\#j}\right\}  ,$ although important in finding a
$\mathbf{c}_{\#}$, is not important for our overall optimization procedure as
it is recovered in the optimization of $\mathbf{\theta}$ and $\mathbf{\psi}$.


\end{document}